KernelBench v3 contains 41 tasks across four levels (Table~\ref{tab:levels}). Each task provides a PyTorch reference model and an input generator. The model under evaluation must produce a CUDA implementation that matches the reference output and improves runtime. Level 1 contains basic kernels such as matmul, softmax, and normalization. Level 2 contains fused multi-op patterns such as matmul-plus-activation-plus-norm chains. Level 3 covers single attention blocks and a transformer block. Level 4 focuses on modern inference kernels, including DeepSeek-style MLA and MoE patterns~\cite{deepseek-v3}, GQA~\cite{gqa}, FP8 and INT4 matmuls~\cite{fp8-formats,int4-quant}, and linear-attention variants such as Gated DeltaNet and Kimi Delta Attention~\cite{gated-deltanet,kimi-linear}.

Some Level 4 baselines are already optimized Triton implementations~\cite{triton}. For example, the Gated DeltaNet and Kimi Delta Attention tasks use optimized kernels rather than naive PyTorch baselines, which sets a high bar for LLM-generated code. As a result, Level 4 performance reflects whether models can compete with production-quality kernels rather than only outperform naive baselines.

\begin{table}[t]
\centering
\caption{KernelBench v3 levels and turn limits.}
\label{tab:levels}
\begin{tabular}{cccc}
\toprule
Level & Tasks & Turns & Description \\
\midrule
L1 & 15 & 10 & Simple ops: matmul, softmax, conv, norms \\
L2 & 15 & 12 & Fused ops: matmul and activation chains \\
L3 & 3 & 15 & Single blocks: attention and transformer block \\
L4 & 8 & 15 & Novel layers: MLA, MoE, GQA, FP8, INT4, DeltaNet \\
\bottomrule
\end{tabular}
\end{table}

KernelBench v3 evaluates on two target GPUs: H100 (Hopper) and B200 (Blackwell)~\cite{nvidia-h100,nvidia-b200}. The benchmark runs in a Modal sandbox with CUDA 13.1~\cite{modal-docs,cuda-13-1}, providing a consistent environment for compilation and benchmarking.
